\twocolumn[{
\section{Results: Outcome Prediction}
\begin{center}\small
\begin{tabular}{lrr}
\toprule
Model / feature representation & Validation equal-phase log loss & Reported public log loss\\
\midrule
Phase-specific word TF-IDF LR baseline & 0.58377 & 0.60149\\
50/50 concatenated-word / separate-field blend & 0.58227 & 0.59962\\
25/25/50 word / field / character blend & 0.58090 & 0.59696\\
50/50 historical blend / phase-shared text model & \textbf{0.56839} & \textbf{0.59011}\\
\bottomrule
\end{tabular}
\captionof{table}{The final outcome model retains the historical text blend as one component.
Lower loss is better. Public values come from user-provided Kaggle screenshots; private results
are unknown. Historical delivery files remain separately reproducible.}
\label{tab:p2models}
\end{center}
\begin{center}\small
\begin{tabular}{llrrrr}
\toprule
Model & Phase & F1 & ROC-AUC & AP & Log loss \\
\midrule
Baseline & I & 0.7795 & 0.7820 & 0.8285 & 0.5966 \\
Baseline & II & 0.6236 & 0.7236 & 0.7088 & 0.6082 \\
Baseline & III & 0.8126 & 0.7663 & 0.8685 & 0.5464 \\
Current phase-shared blend & I & 0.7742 & 0.7961 & 0.8400 & 0.5738 \\
Current phase-shared blend & II & 0.6297 & 0.7387 & 0.7259 & 0.5983 \\
Current phase-shared blend & III & 0.8178 & 0.7817 & 0.8776 & 0.5330 \\
\bottomrule
\end{tabular}

\captionof{table}{Required validation metrics by phase for the original baseline and current
public-best outcome model. Success is the positive class; F1 uses 0.5, ROC-AUC and average
precision (AP) use probabilities. Validation has 149/501/363 phase-I/II/III rows.}
\label{tab:p2phase}
\end{center}
}]
\subsection{Models and feature choices}
Part 2 has 8,037 training and 1,013 validation trial-phase rows. All phase copies of a trial stay
together in model-selection splits. Test has 1,500 rows, 500 per phase, so equal-phase mean
log loss is the selection objective. Observed test start years are 2014--2021, later than the
training/validation population through 2013.

The baseline fits one TF-IDF LR per phase on concatenated eligibility criteria, disease, and
drug text~\cite{tfidf,sklearn}. Word unigrams/bigrams use minimum document frequency 3,
at most 20,000 features, lowercase Unicode accent normalization, sublinear TF, smoothed IDF,
and L2 normalization. Missing text is empty. LR uses $C=1$, liblinear, 500 iterations, and
tolerance $10^{-4}$. Reason labels and reason embeddings are excluded from outcome predictors.

Separate-field models use criteria/disease/drug vocabulary limits 20,000/5,000/5,000 and block
scaling $1/\sqrt{3}$. Character models use word-boundary 3--5 grams, minimum document frequency 3,
and at most 30,000 features per phase. The historical blend weights concatenated word,
separate-field word, and character probabilities 0.25/0.25/0.50. Text plus metadata and pathway
features did not improve validation loss: the $C=1$ metadata model scored 0.59460, compared with
0.58377 for text alone. Exact drug matching and sparse pathway coverage can limit added signal;
these are plausible explanations rather than isolated causal effects.

\subsection{Phase sharing and later cohorts}
The new text LR component shares coefficients across phases, adds phase-specific interactions
scaled by 0.35 and phase intercepts, and uses training recency weights with a six-year half-life.
Start year changes training weights, not inference features. This permits common text effects
while retaining phase differences. Its final probability receives weight 0.5 alongside the
historical blend.

A historical screen trains through 2009 and evaluates 2010--2013; equal-phase loss improves
from 0.58425 to 0.57288 for the selected blend. Five trial-grouped training folds select the
mixture; OOF loss improves from 0.59649 to 0.59155. The supplied validation check then improves
from 0.58090 to 0.56839. These comparisons motivated the public submission, which scores
0.59011. Outcome training does not add the supplied validation labels.

\subsection{Overall performance and generalization}
Across all 1,013 validation trial-phase rows, pooled F1/ROC-AUC/AP improve from
0.7354/0.7569/0.7982 to 0.7387/0.7737/0.8144. These pooled metrics differ from the
equal-phase mean log loss used for selection. Lower loss need not improve every phase's F1.

Public log loss is 0.59011 versus validation 0.56839, a gap of 0.02172.
Later trial cohorts, repeated validation selection, and public-subset sampling variability
may contribute; this comparison does not isolate any one cause. Private performance is unknown.
